% Generated draft: requires mathematical and source review.
\documentclass[11pt]{article}
\usepackage[margin=25mm]{geometry}
\usepackage{fontspec}
\setmainfont{Latin Modern Roman}
\usepackage{amsmath,amssymb,mathtools,unicode-math}
\setmathfont{Latin Modern Math}
\usepackage{enumitem,needspace,xurl,hyperref}
\setlength{\parindent}{0pt}
\setlength{\parskip}{5pt}
\setlength{\emergencystretch}{3em}
\newcommand{\N}{\mathbb{N}}
\newcommand{\Z}{\mathbb{Z}}
\newcommand{\Q}{\mathbb{Q}}
\newcommand{\R}{\mathbb{R}}
\newcommand{\C}{\mathbb{C}}
\newcommand{\F}{\mathbb{F}}
\newcommand{\A}{\mathbb{A}}
\newcommand{\PP}{\mathbb{P}}
\newcommand{\E}{\mathbb{E}}
\newcommand{\eps}{\varepsilon}
\newcommand{\dd}{\,\mathrm{d}}
\begin{document}
\section*{UnsolvedMath ID: 30001517}
\textbf{Source problem code:} catalogue entry
\subsection{Definitions and mathematical statement}
The supplied record asks about bodies invisible or of zero resistance in several directions, but does not specify the ambient dimension, the admissible class of bodies, or the scattering and measure conventions. Consequently, the following is a conditional formulation using the usual exterior specular-billiard interpretation; these conventions require confirmation against the source.

Let \(n\geq 2\) be an unspecified ambient dimension, let \(S^{n-1}=\{v\in\R^n:\|v\|=1\}\), and let \(\mathcal B_n\) denote the intended, presently unspecified class of bodies in \(\R^n\). For \(v\in S^{n-1}\), parametrize incoming oriented lines by \(b\in v^\perp\), with incoming motion \(x(t)=b+tv\). At a regular boundary collision, the specular reflection law is
\[
 u_{\mathrm{out}}=u_{\mathrm{in}}-2\langle u_{\mathrm{in}},\nu\rangle\nu,
\]
where \(\nu\) is a unit normal. Call scattering in direction \(v\) admissible when, for Lebesgue-almost every \(b\in v^\perp\), the trajectory is well defined and eventually follows an outgoing straight ray. Write its outgoing unit velocity as \(v_B^+(b,v)\).

Under this interpretation, \(B\) has \emph{zero resistance in direction \(v\)} when scattering is admissible and
\[
 v_B^+(b,v)=v\qquad\text{for Lebesgue-almost every }b\in v^\perp.
\]
It is \emph{invisible in direction \(v\)} when, in addition, the outgoing ray lies on the original oriented line \(b+\R v\) for almost every \(b\). Thus invisibility permits a travel-time delay but requires restoration of both direction and line. These are candidate scattering definitions, not definitions explicitly supplied by the record; in particular, their applicability to the intended boundary regularity and to exceptional or trapped trajectories must be checked.

Define the direction sets
\[
 I(B)=\{v\in S^{n-1}:B\text{ is invisible in direction }v\},\qquad
 Z(B)=\{v\in S^{n-1}:B\text{ has zero resistance in direction }v\}.
\]
With these conventions, the questions are:
\begin{enumerate}
\item Does there exist \(B\in\mathcal B_n\) with \(\#I(B)\geq 3\)? Does there exist, possibly a different body, \(B\in\mathcal B_n\) with \(\#Z(B)\geq 3\)?
\item Is it true that every \(B\in\mathcal B_n\) satisfies
\[
 \sigma_{n-1}(I(B))=0
 \qquad\text{and}\qquad
 \sigma_{n-1}(Z(B))=0,
\]
where \(\sigma_{n-1}\) is spherical surface measure, provided these direction sets are measurable?
\end{enumerate}
The source does not specify whether directions are oriented or identified under \(v\sim -v\), whether spherical surface measure is intended, or whether outer measure should replace measure for potentially nonmeasurable direction sets. The dimension, quantification over dimensions, and class \(\mathcal B_n\) remain unresolved. The supplied background specifically distinguishes the original smooth-body invisibility question from later fractal constructions and from a two-dimensional piecewise-smooth result. Those classes must not be interchanged.
\subsection{Short English statement}
The problem asks whether a body can be invisible, or have zero resistance, in at least three directions, and whether either property must occur only on a set of directions of measure zero. The supplied background refers to an original smooth-body invisibility question and reports partial progress in other classes. A definitive statement still requires the intended dimension, body regularity, direction convention, and meaning of positive measure.
\subsection{Sources}
\begin{itemize}
\item Supplied local UnsolvedMath registry record, permanent ID \(30001517\), title ``Multidirectional Invisible and Zero Resistance Bodies,'' source problem code \texttt{catalogue entry}, in the Oberwolfach Reports Open Problems collection. The record gives proposed year \(2010\) and names no proposer. Its statement is labeled ``exact,'' with the supplied verification note: ``The statement is internally readable; retained byte-for-byte.'' This is a registry qualification, not independent source verification.
\item Supplied report citation: \emph{Calculus of Variations (2011)}, Oberwolfach Reports, DOI \texttt{10.4171/owr/2010/31}, pp.\ 1827--1884. \url{https://doi.org/10.4171/owr/2010/31}. The background describes Oberwolfach Reports as documenting mathematical workshops and their open-problem discussions. The supplied year \(2011\), proposed year \(2010\), and DOI containing \(2010\) are retained without reconciliation.
\item Plakhov--Roshchina, \emph{Invisibility in billiards}. \url{https://arxiv.org/abs/1009.1104}.
\item Plakhov--Roshchina, \emph{Fractal bodies invisible in 2 and 3 directions}. \url{https://arxiv.org/abs/1107.5667}.
\item Plakhov--Roshchina, \emph{Invisibility in billiards is impossible in an infinite number of directions}. \url{https://arxiv.org/abs/1206.6965}.
\item Literature assessment dated exactly \texttt{2026-08-21}: the supplied record labels the problem ``partially solved'' and reports that the original smooth-body invisibility question in three or more directions remains unsettled in that class, that fractal bodies invisible in three orthogonal directions have been constructed, and that finiteness has been proved in a two-dimensional piecewise-smooth setting. This is supplied dated curation derived from the OWR record, not an independently verified assessment. The record expressly calls for independent verification of time-sensitive status and cited claims; it does not provide a corresponding detailed assessment of every zero-resistance question.
\end{itemize}
\end{document}
